% GDZ PPN236018647,native577lower–578;HKUST578–579.
\NumberedItem{108}
\indent Numeratorem et denominatorem sinus lemniscatici (universalissime accepti) ad quantitates integrales reducere contigit; simul omnium functionum lemniscaticarum, quae excogitari possunt, evolutiones in series infinitas ex principiis genuinis haustae; inventum pulcherrimum sane nullique praecedentium inferius.
\par\indent Praeterea iisdem diebus principia deteximus, secundum quae series arithmetico-geometricae interpolari debent, ita ut terminos in progressione data ad indicem quemcumque rationalem pertinentes per aequationes algebraicas exhibere iam in potestate sit.
\par\noindent\hfill[1800] Mai. ult. Iun. 2. 3.

\medskip\begingroup\small
\indent Der \textit{sinus lemniscaticus universalissime acceptus} ist die in der Scheda Ac, S. 26, siehe oben S. 194, nach der Analogie des sinus lemniscaticus aufgebaute Funktion \(S\varphi\). Die Zurückführung des Zählers und
\par\endgroup

\clearpage\SourceLeaf{578}\ArithmeticHead{549}{R}{1800 MAI. 16—1800 MAI. ULT. IUN. 2. 3.}
\begingroup\small
\noindent des Nenners dieser Funktion, d. h. der Funktionen \(T\varphi\) und \(W\varphi\), auf \textit{quantitates integrales} liefert den unmittelbaren Nachweis dafür, daß ihr Quotient \(x=S\varphi\) der Differentialgleichung
\[
\left(\frac{dS\varphi}{d\varphi}\right)^2=(1+\mu^2x^2)(1-x^2)
\]
\noindent Genüge leistet. Dieser Nachweis ist in der Scheda Ac nicht aufgezeichnet und hat sich auch in keinem andern aus dieser Zeit herrührenden Teile des Nachlasses gefunden. Er ist uns jedoch in zwei späteren Fassungen erhalten, die sich beide in dem Handbuch 16, Bb (begonnen November 1801) befinden. Die ältere, die wahrscheinlich aus dem Jahre 1825 stammt, steht S. 111, 112 des Handbuchs und ist mit Änderungen gegen die Handschrift Werke III, S. 401, 402, unverändert oben S. 308—310 abgedruckt; die spätere, die nach \textsc{Gauss}’ eigener Angabe im August 1827 verfaßt ist, steht S. 139 des Handbuchs und ist Werke III, S. 473, art. [5.] abgedruckt. Daß wir aber diesen Nachweis für die Zeit unserer Tagebuchnotiz in Anspruch nehmen dürfen, folgt daraus, daß er wesentlich auf den Formeln für die Transformation zweiter Ordnung der Funktionen
\[
T\varphi,\ W\varphi,\ T\left(\frac\varpi2-\varphi\right),\ W\left(\frac\varpi2-\varphi\right)
\]
\noindent beruht, die auf S. 31 der Scheda Ac aufgezeichnet sind (siehe oben S. 195, 196, art. [3.]), so daß dieser Umstand erst den Grund dafür erkennen läßt, weshalb \textsc{Gauss} diese Transformationsformeln in der Scheda Ac entwickelt. Ferner wird die für die Ausführung jenes Nachweises erforderliche Formel, S. 111 des Handbuchs 16, Bb,
\[
ab'-ba'=\frac12ab(a^4-b^4),
\]
\noindent (siehe oben S. 310) von der \textsc{Gauss} a. a. O. sagt, daß \textls[120]{der Beweis davon tiefer liegt,} in der Scheda Af, S. 13 (abgedruckt oben S. 212, art. [5.]) als besonderes \textit{Theorema} abgeleitet; diese Scheda Af ist aber ganz im Jahre 1801 geschrieben worden. —
\par\indent Die im zweiten Satze des ersten Absatzes unseres Tagebuchtextes bezeichnete Entwicklung »aller nur denkbaren lemniskatischen Funktionen in unendliche Reihen aus genuinen Prinzipien abgeleitet« bezieht sich auf die S. 41 und 43 der Scheda Ac befindlichen Rechnungen (abgedruckt oben S. 202, art. [10.] und S. 204, art. [12.]), die auf die Ableitung der berühmten Identität zwischen den Reihen- und den Produktentwicklungen der Thetafunktionen:
\[
\begin{aligned}
&(1+\alpha x)(1+\alpha x^3)(1+\alpha x^5)\ldots\left(1+\frac x\alpha\right)\left(1+\frac{x^3}\alpha\right)\left(1+\frac{x^5}\alpha\right)\\
&\qquad=P\left(1+x\left(\alpha+\frac1\alpha\right)+x^4\left(\alpha^2+\frac1{\alpha^2}\right)+x^9\left(\alpha^3+\frac1{\alpha^3}\right)+\ldots\right),
\end{aligned}
\]
\noindent wo \(P\) von \(\alpha\) unabhängig ist, hinzielen. Vermöge dieser Identität ist \textsc{Gauss} nämlich imstande, für eine beliebige Thetafunktion die durch Kenntnis ihrer Nullstellen gegebene Produktentwicklung in eine Reihenentwicklung umzusetzen.
\par\indent Der zweite Absatz der Notiz Nr. 108 betrifft die Gleichung für die Teilung der Perioden. Eine darauf bezügliche Aufzeichnung aus dieser Zeit ist uns nicht erhalten.
\par\noindent\hfill\textsc{Klein.\quad Schlesinger.}
\par\endgroup
